% TOP_PROBLEM: {"id": 1162, "title": "Hilbert's Sixteenth Problem", "category_id": 4, "category": {"id": 4, "name": "algebra", "display_name": "Algebra", "description": "Group theory, ring theory, field theory, and algebraic structures.", "slug": "algebra", "order_index": 4, "created_at": "2026-07-31T15:26:25.671Z"}, "status": "open"}
% FIRST_POSTED: 2026-10-06
% ATTEMPT_STATUS: unresolved
% OPENAI_CLAIM: {"schema_version":1,"source_repo":"https://github.com/openai/math","source_commit":"adc7f1241b42e322a6451854ab7e4b4c146bf78a","release_date":"2026-10-06","match":"partial","resolution":"partial","independently_reviewed":false,"second_pass":"n/a","families":[{"family":"143","title":"Hilbert's sixteenth problem: uniform bounds for limit cycles","subject":"Dynamical systems and ergodic theory","manuscripts":[{"title":"Two limit cycles for quintic Liénard systems","pdf_path":"preprints/two-limit-cycles-for-quintic-lienard-systems-September-24-2026/two-limit-cycles-for-quintic-lienard-systems-September-24-2026.pdf","readme_path":"preprints/two-limit-cycles-for-quintic-lienard-systems-September-24-2026/README.md","date":"2026-09-24","theorem_ref":"thm:main (sections/introduction.tex:13)","statement_sources":[{"path":"preprints/two-limit-cycles-for-quintic-lienard-systems-September-24-2026/build/sections/introduction.tex","line":13,"label":"thm:main"}]},{"title":"Uniform bounds for planar polynomial limit cycles","pdf_path":"preprints/uniform-bounds-for-planar-polynomial-limit-cycles-September-24-2026/uniform-bounds-for-planar-polynomial-limit-cycles-September-24-2026.pdf","readme_path":"preprints/uniform-bounds-for-planar-polynomial-limit-cycles-September-24-2026/README.md","date":"2026-09-24","theorem_ref":"thm:uniform (sections/01-introduction.tex:12)","statement_sources":[{"path":"preprints/uniform-bounds-for-planar-polynomial-limit-cycles-September-24-2026/build/sections/01-introduction.tex","line":12,"label":"thm:uniform"}]}],"lean":{"doc_path":"lean/docs/143.md","comparators":["lean/ComparatorChallenges/QuinticLienard.json","lean/ComparatorChallenges/QuinticLienard.lean"],"declarations":["OAI.QuinticLienard.main"],"covers_main_theorem":false},"reasoning_trace":null}]}
\documentclass[11pt]{article}
\usepackage[margin=25mm]{geometry}
\usepackage{fontspec}
\setmainfont{Latin Modern Roman}
\usepackage{amsmath,amssymb,mathtools}
\usepackage{unicode-math}
\setmathfont{Latin Modern Math}
\usepackage{enumitem}
\usepackage{xurl,hyperref,bookmark}
\hypersetup{colorlinks=true,urlcolor=blue}
\setcounter{secnumdepth}{0}
\setlength{\parindent}{0pt}
\setlength{\parskip}{5pt}
\setlength{\emergencystretch}{3em}
\begin{document}
\section*{Hilbert's Sixteenth Problem}

\subsection{OpenAI's claim}
Manuscript 1 (family 143, 2026-09-24): A classical Liénard system with arbitrary real primitive polynomial of degree at most five is claimed to have at most two isolated periodic orbits, and two are claimed to occur. Manuscript 2 (family 143, 2026-09-24): For each polynomial degree, a finite bound on the number of isolated periodic orbits of all real planar polynomial vector fields of that degree is claimed.

\subsection{Mathematical statement of the claimed result}
\paragraph{Mathematical claim 1.} For any real polynomial $F$ of degree at most five, consider
\[
\dot x=y-F(x),\qquad\dot y=-x.
\]
OpenAI claims at most two limit cycles in $\mathbb R^2$, with examples attaining two. A cycle means the geometric image of a nonconstant periodic trajectory isolated among periodic orbits; multiplicity and stability do not affect the count. Periodic trajectories of a center are excluded.

\paragraph{Mathematical claim 2.} For each integer $d\ge1$, OpenAI asserts existence of a finite nonnegative integer $B(d)$ satisfying
\[
\#\{\text{isolated periodic-orbit images of }\dot x=P(x,y),\ \dot y=Q(x,y)\}\le B(d)
\]
for every pair of real polynomials $P,Q$ of total degree at most $d$. The count covers the whole plane and ignores multiplicity. The theorem is a qualitative uniform bound; it does not specify a sharp numerical value of $B(d)$.

\subsection{Outline of the argument}
Manuscript 1 (family 143, 2026-09-24): A quadratic midpoint model tracks fitted slope and curvature by transport equations. Endpoint variation, Riccati linearization, and a Schwarzian identity control the width dependence of the model coefficients, allowing the two-cycle argument without the odd-polynomial restriction of the older theorem.

Manuscript 2 (family 143, 2026-09-24): Nested complex domains separate the required asymptotic scales, while high real jets retain a fixed coherence order through rapidly rotating dependent scales. Proper barriers and generic tilts reduce parameters in a derivative-stable analytic class; finite-dimensional counting then turns absolute finiteness into a uniform degree bound.

\subsection{Scope relative to this problem}
The quintic Li\'enard paper determines a maximum of two cycles for its specified subclass. The other selected manuscript asserts a finite bound for all planar polynomial vector fields at each fixed degree. Neither claim determines the sharp degree-by-degree maximum or the possible arrangements of cycles asked for here.

\subsection{Formal verification}
For family 143, the scope page documents the following selections. The selection concerns degree-at-most-five Liénard systems: two isolated cycles are the maximum and that maximum occurs. It is not the arbitrary-degree planar polynomial bound. The documented statement covers the main claim of manuscript 1. Its selected declaration also covers the clause used for this problem. Manuscript 2 is not among the papers selected by that page, so the family link does not formalize its main conclusion. It does not establish formal coverage of the complete problem clause. Comparator files: QuinticLienard.lean. Declarations: OAI.QuinticLienard.main.

This site has not executed the Lean files or checked their dependencies. These descriptions report the scope of the linked declarations, rather than a new verification of the proofs.

\subsection{Status caveat}
OpenAI attributes these results to an internal model. Its release notice cautions that results lacking full formalization can contain mistakes. This site has not independently assessed the proofs; a solved label records the scope of the reported claim.

\subsection{Sources}
\begin{itemize}
\item OpenAI, \emph{Two limit cycles for quintic Liénard systems}, 2026-09-24; citation key \texttt{OAI:two-limit-cycles-for-quintic-lienard-systems-September-24-2026}. \href{https://github.com/openai/math/blob/adc7f1241b42e322a6451854ab7e4b4c146bf78a/preprints/two-limit-cycles-for-quintic-lienard-systems-September-24-2026/two-limit-cycles-for-quintic-lienard-systems-September-24-2026.pdf}{PDF (GitHub)}; \href{https://github.com/openai/math/raw/adc7f1241b42e322a6451854ab7e4b4c146bf78a/preprints/two-limit-cycles-for-quintic-lienard-systems-September-24-2026/two-limit-cycles-for-quintic-lienard-systems-September-24-2026.pdf}{PDF (direct)}; \href{https://github.com/openai/math/blob/adc7f1241b42e322a6451854ab7e4b4c146bf78a/preprints/two-limit-cycles-for-quintic-lienard-systems-September-24-2026/README.md}{Manuscript page and citation}; \href{https://github.com/openai/math/blob/main/preprints/two-limit-cycles-for-quintic-lienard-systems-September-24-2026/two-limit-cycles-for-quintic-lienard-systems-September-24-2026.pdf}{Latest version}; Statement source: \href{https://github.com/openai/math/blob/adc7f1241b42e322a6451854ab7e4b4c146bf78a/preprints/two-limit-cycles-for-quintic-lienard-systems-September-24-2026/build/sections/introduction.tex\#L13}{thm:main}.
\item OpenAI, \emph{Uniform bounds for planar polynomial limit cycles}, 2026-09-24; citation key \texttt{OAI:uniform-bounds-for-planar-polynomial-limit-cycles-September-24-2026}. \href{https://github.com/openai/math/blob/adc7f1241b42e322a6451854ab7e4b4c146bf78a/preprints/uniform-bounds-for-planar-polynomial-limit-cycles-September-24-2026/uniform-bounds-for-planar-polynomial-limit-cycles-September-24-2026.pdf}{PDF (GitHub)}; \href{https://github.com/openai/math/raw/adc7f1241b42e322a6451854ab7e4b4c146bf78a/preprints/uniform-bounds-for-planar-polynomial-limit-cycles-September-24-2026/uniform-bounds-for-planar-polynomial-limit-cycles-September-24-2026.pdf}{PDF (direct)}; \href{https://github.com/openai/math/blob/adc7f1241b42e322a6451854ab7e4b4c146bf78a/preprints/uniform-bounds-for-planar-polynomial-limit-cycles-September-24-2026/README.md}{Manuscript page and citation}; \href{https://github.com/openai/math/blob/main/preprints/uniform-bounds-for-planar-polynomial-limit-cycles-September-24-2026/uniform-bounds-for-planar-polynomial-limit-cycles-September-24-2026.pdf}{Latest version}; Statement source: \href{https://github.com/openai/math/blob/adc7f1241b42e322a6451854ab7e4b4c146bf78a/preprints/uniform-bounds-for-planar-polynomial-limit-cycles-September-24-2026/build/sections/01-introduction.tex\#L12}{thm:uniform}.
\item Family 143: \href{https://github.com/openai/math/blob/adc7f1241b42e322a6451854ab7e4b4c146bf78a/lean/docs/143.md}{Lean scope}.
\item \href{https://github.com/openai/math/blob/adc7f1241b42e322a6451854ab7e4b4c146bf78a/lean/ComparatorChallenges/QuinticLienard.json}{Comparator QuinticLienard.json}.
\item \href{https://github.com/openai/math/blob/adc7f1241b42e322a6451854ab7e4b4c146bf78a/lean/ComparatorChallenges/QuinticLienard.lean}{Comparator QuinticLienard.lean}.
\item \href{https://github.com/openai/math/blob/adc7f1241b42e322a6451854ab7e4b4c146bf78a/CONTENTS.md}{OpenAI catalogue entry}.
\end{itemize}
\end{document}
